% =====================================================================
%  INTERNAL REVIEWER CHECKLIST -- NOT FOR SUBMISSION.
%  Rendered only when \showreviewtrue is set in main.tex.
%  Mark an issue addressed by adding \done{<ID>} below, then recompile.
% =====================================================================
% \done{A5}
% \done{A8}
% ---------------------------------------------------------------------
\clearpage
\section*{\color{openred}\rule{\linewidth}{1pt}\\[2pt]Internal Reviewer Checklist --- remove before submission\\[2pt]\rule{\linewidth}{1pt}}
\addcontentsline{toc}{section}{Internal Reviewer Checklist (remove before submission)}

\noindent Severity is encoded in the section ordering: Section~\ref{sec:critical} items are blocking (internal contradictions a reviewer will catch immediately); Sections~\ref{sec:stats}--\ref{sec:claims} affect the soundness of the claims (TMLR's central acceptance criterion is that claims are supported by evidence); the remainder are polish, compliance, and clarity.

\section{Critical internal inconsistencies (blocking)}
\label{sec:critical}
\begin{issuelist}

\issue{A1}{\textbf{Claude Sonnet's HLE status contradicts itself across the paper.} Five locations treat Sonnet as having HLE data with a measurable slope of $+6.11^\dagger$: \texttt{results.tex:76} (``the five remaining models''), \texttt{limitations.tex:7}, \texttt{tab:rq4\_hle}, \texttt{methods.tex:249} (which excludes only Haiku), and \texttt{tab:hle\_consensus} ($n{=}61$ $T_0$-correct queries). Only two mark it $\varnothing$: the Sonnet row of Table \texttt{tab:rq3} (``no $T_0$-correct questions'') and \texttt{tables/coverage.tex}. The smoking gun is that \texttt{tab:rq3}'s own caption names \emph{only} Haiku as $\varnothing$, directly contradicting its own Sonnet body row; the stale $\varnothing$ traces to a commented-out ancestor in \texttt{limitations.tex:3} that grouped Sonnet with Haiku as excluded.}{The five-location majority is correct: Sonnet has HLE data. Set the \texttt{tab:rq3} Sonnet HLE row to ``$+6.11^\dagger$ / Capitulator$^\dagger$'' (matching \texttt{results.tex:76}) and flip the \texttt{coverage.tex} Sonnet--HLE cell to $\checkmark$. Then add a single source-of-truth coverage macro or table that every section cites, so the numbers cannot diverge again.}

\issue{A2}{\textbf{Claude Haiku's HLE exclusion is glossed with the wrong reason.} The $\varnothing$ symbol for Haiku's HLE cell is explained as ``no $T_0$-correct questions'' (\texttt{tab:rq3}) / ``essentially no $T_0$-correct questions'' (\texttt{methods.tex}), but \texttt{results.tex} and \texttt{tab:hle\_consensus} show Haiku has $33$ $T_0$-correct queries --- it loses consensus on \emph{all} of them by $\approx T_2$, which is what actually makes the trajectory unmeasurable. \texttt{methods.tex:249} compounds this with a false causal claim --- that Haiku ``cannot answer HLE questions well enough at baseline to have a correct-at-$T_0$ subset'' --- directly contradicted by $n{=}33$. (The same $\varnothing$-mislabel also affects Sonnet; see A1.)}{Use one canonical sentence wherever the exclusion is explained: ``Haiku has $33$ $T_0$-correct HLE queries but loses consensus on all of them by $\sim T_2$, so no per-turn flip trajectory is measurable.'' Reword the \texttt{tab:rq3} footnote and the \texttt{coverage.tex} caption so $\varnothing$ means ``trajectory unmeasurable due to consensus loss,'' not ``no $T_0$-correct questions,'' and delete the false baseline-accuracy clause at \texttt{methods.tex:249}.}

\issue{A3}{\textbf{RQ1 ``no certain group'' claim is false as stated, and the GPT-5.4 cell is likely a data error.} \texttt{results.tex} (RQ1) states that on GPQA-Diamond ``Claude Haiku and GPT-5.4 have no certain group (every question elicits non-zero baseline entropy),'' with Table \texttt{tab:rq1} showing N/A for both --- but Table \texttt{tab:entropy\_dist} reports Haiku with $22$ certain questions ($27.8\%$) and GPT-5.4 with $150$ certain ($75.8\%$). The stated reason (``every question is uncertain'') is therefore flatly false. The two cells then diverge: Haiku's N/A \emph{may} be legitimate, since ``certain'' ($H{=}0$) can mean certain-\emph{but-wrong}, and a weak model on GPQA could have all 22 certain questions be wrong, leaving an empty certain-\emph{correct} group. But GPT-5.4's N/A is almost certainly a pipeline error: for it to be real, the strongest model would have to be confidently wrong on all 150 certain questions, while Sonnet ($33$ certain), Mini ($139$), and Nano ($155$) all populate their certain-correct cells, and GPT-5.4 \emph{passes} the paper's own eligibility gate ($150$ vs.\ $48$, both non-empty) so it should not be N/A at all.}{Re-run the GPT-5.4 GPQA certain-correct cell and report the real flip rate (this is a likely-wrong number in a headline RQ1 cell). Reword the false justification in \texttt{results.tex} and the \texttt{tab:rq1} caption: N/A must mean ``no question is both certain ($H{=}0$) and correct-at-$T_0$,'' not ``every question is uncertain.'' Add a certain-\emph{correct}-at-$T_0$ count column to \texttt{tab:entropy\_dist} so each N/A can be checked directly.}

\issue{A4}{\textbf{Sample sizes for the same cell disagree between RQ1 and RQ2 by $\sim$22$\times$, and the RQ1 unit-of-analysis is self-contradictory.} For Claude Haiku / MMLU-Pro, \texttt{tab:rq2\_combined} reports $n=294$ turn-level observations ($\approx$10 questions at $R{=}5\times6$ turns $=30$ per question), but the RQ1 uncertain Wilson CI $28.2\%$ $[27.1, 29.3]$ (half-width $1.1$\,pp) back-solves to $N\approx6{,}400$ pooled turn-level observations ($\approx209$ questions $\times30$); the certain-group CI $[0.0,0.1]$ implies thousands as well. A genuine per-question CI ($N{=}209$) would be $\approx[22,34]$, six times wider --- so RQ1's CIs are pooled at the \emph{turn} level, which contradicts \texttt{experiments.tex:6} (``the question is the unit of analysis'') while matching \texttt{methods.tex:208} (``Wilson CIs\ldots treat these observations as independent''). Every other RQ2 MMLU-Pro cell divides cleanly by $30$ into a plausible question count; only Haiku collapses to $\sim$10, so under the identical claimed correct-at-$T_0$ restriction either RQ2's $n$ is wrong or the filter was applied to only one of the two analyses.}{Audit both pipelines against the authors' code (not resolvable from the manuscript). Report the per-cell $n$ for RQ1 in \texttt{tab:rq1} alongside the RQ2 $n$ so the two denominators are checkable, and state the exact unit of analysis for every reported statistic --- in particular reconcile \texttt{experiments.tex:6} with \texttt{methods.tex:208}. If the Wilson CIs are pooled at the turn level, recompute them per-question (or cluster-robust by question); the intervals will widen substantially, consistent with the authors' own optimistic-SE caveat.}

\issue{A5}{\textbf{HLE sample-size inconsistency exists only in dead text (downgraded).} A divergent range ``28--55'' appears in the paper, but only inside the commented-out draft at \texttt{limitations.tex:1--3}, which never compiles. All three \emph{live} occurrences (\texttt{results.tex:76}, \texttt{discussion.tex:4}, \texttt{limitations.tex:7}) agree on ``24--72,'' which matches the true min--max of \texttt{tab:hle\_consensus} (per-model $n$ = 33/61/72/24/26/28; min $=24$ Nano, max $=72$ Gemini). No live-text inconsistency remains; this collapses into D4 (remove commented-out drafts).}{Delete the commented-out draft block at \texttt{limitations.tex:1--3} (also resolves part of D4). Optionally, cross-reference \texttt{tab:hle\_consensus} instead of hard-coding ``24--72'' in three places, so the range cannot drift if the table changes.}

\issue{A6}{\textbf{Eligibility-criterion wording differs between methods and the coverage caption (cosmetic).} \texttt{methods.tex} defines eligibility as two-group non-emptiness ($H=0$ and $H<0$ both non-empty) and explicitly says the five quantile bins do \emph{not} enter the test, yet the coverage caption defines $\ast$ as ``single populated entropy bin.'' For the actual $\ast$ cells (Gemini, baseline unanimous at $H{=}0$) the certain group and the lowest bin are the same set, so the exclusion outcome is correct --- only the terminology is imprecise. Low severity, no result affected.}{Rewrite the \texttt{coverage.tex} caption to match the methods definition: ``single populated group ($H{=}0$ only or $H{<}0$ only).''}

\issue{A7}{\textbf{Consensus loss / $B_t$ is defined on letter-only dialogues but applied to (and equated with) CoT dialogues.} \S Multi-Turn Pressure Protocol defines $B_t$ and consensus loss over the $R=5$ letter-only re-poll dialogues, yet the HLE fragmentation paragraph in \texttt{results.tex} computes consensus loss ``across the $R = 5$ \emph{CoT} dialogues'' and \texttt{tab:hle\_consensus} says ``five repeated CoT samples.'' The mismatch also lives inside Methods itself: \texttt{methods.tex:143--144} says the letter-only $B_t$ ``is the Bel.\ Ent.\ column'' of the RQ4 tables, but those tables and \texttt{results.tex:91} describe that column as belief entropy over the CoT samples. The two run sets have different inner aggregation ($K_{\text{turn}}{=}5$ re-polls vs.\ one FINAL ANSWER per CoT trace), so a single unqualified definition cannot cover both.}{Parameterize the Methods definition of $B_t$/consensus loss by run set: the dialogue-level belief is the modal answer of the $K_{\text{turn}}{=}5$ re-polls for letter-only runs, and the FINAL ANSWER of the trace for CoT runs. Fix \texttt{methods.tex:143--144} so the RQ4 Bel.\ Ent.\ column is attributed to the CoT ensemble, and keep the explicit ``CoT'' labels in Results/appendix as reinforcement.}

\issue{A8}{\textbf{Duplicate bibliography entry.} \texttt{bib.bib} contains \texttt{jiang2025feedbackfriction} and \texttt{jiang2025feedbackfrictionllmsstruggle}, both pointing to arXiv:2506.11930 with the same title and authors. Only \texttt{jiang2025feedbackfriction} is cited in the text.}{Delete the uncited \texttt{jiang2025feedbackfrictionllmsstruggle} entry and keep the single cited key.}

\end{issuelist}

\section{Statistical methodology}
\label{sec:stats}
\begin{issuelist}

\issue{B1}{\textbf{Non-independence is acknowledged but not corrected.} All headline inference (Wilson CIs, per-model logistic SEs, Mann--Whitney $U$) pools turn-level outcomes that are strongly correlated within dialogues and within questions. The paper concedes the SEs are ``likely optimistic'' and defers the fix to future work --- but the fix (cluster-robust SEs by question; question-level random intercept) requires no new data.}{Rerun the RQ1 contrast and RQ2 regressions with by-question cluster-robust standard errors (and add a question-level random effect to the GLMM). Report these as the main numbers; the current naive SEs can move to an appendix. This is the single most likely reviewer request.}

\issue{B2}{\textbf{The pooled-interaction significance story is fragile and inconsistent across sections.} The GLMM fixed-slope gives $\hat\beta_{PH}=0.010$, $p<0.001$; the DerSimonian--Laird pooled estimate gives $p=0.08$; excluding Claude Haiku gives $p=0.005$. The abstract says the pooled interaction ``does not reach conventional significance,'' while the Discussion leans on the post-exclusion $p=0.005$.}{Pick one pre-specified estimator for the headline (the random-effects pooled estimate including all cells, $p=0.08$, is the honest choice), present the Haiku-excluded number strictly as a sensitivity analysis with the exclusion rule (numerical separation) stated \emph{prospectively} in Methods, and make the abstract, Results, and Discussion use identical language.}

\issue{B3}{\textbf{Separation-inflated MLEs are reported rather than fixed.} Claude Haiku MMLU-Pro $\hat\beta_H=-32.5$ and Claude Sonnet GPQA $\hat\beta_H=-12.2$ are acknowledged as divergent MLEs (``direction and significance are valid; magnitude is inflated''), and one cell is dropped entirely.}{Refit all cells with Firth-penalized logistic regression (or weakly-informative Bayesian priors). This yields finite, comparable estimates for every cell, removes the need to drop GPT-5.4 Mini/GPQA, and makes the meta-analysis in B2 cleaner.}

\issue{B4}{\textbf{The certain/uncertain dichotomy is sensitive to $K=25$ and measured with error.} A question where the model's true modal-answer probability is $0.9$ is classified ``certain'' ($H=0$) with probability $0.9^{25}\approx 7\%$ --- so the certain group is contaminated by high-but-not-perfect-confidence questions and vice versa, attenuating or biasing the RQ1 contrast. The justification given for $K=25$ (resolution of $0.04$) does not address classification error at the $H=0$ boundary.}{Add a robustness analysis subsampling $K'\in\{10,15,20\}$ from the existing 25 draws (no new API calls needed) showing the RQ1 gap is stable; and/or treat uncertainty continuously ($H$ or $\hat h$) in the headline analysis rather than as a binary split.}

\issue{B5}{\textbf{The flip metric feeding each analysis is ambiguous.} \S Flip Definition defines an \emph{event-based} flip ($\hat a_{t-1}=a^\ast \wedge \hat a_t = a^\dagger$), then says the RQ3 rate is \emph{state-based}. RQ1's ``per-question flip rate across pressure turns $T_1$--$T_6$'' does not say which; the ``twice the rate'' headline (46.6\% vs.\ 23.6\%) is also an \emph{unweighted average of per-model rates} across heterogeneous models.}{State explicitly, in the Experiments paragraph for each RQ, whether the outcome is event- or state-based. For the headline contrast, report the per-model range alongside a pooled estimate from a model with a random intercept for target model (consistent with B1), rather than an unweighted mean.}

\issue{B6}{\textbf{One-sided Mann--Whitney tests presuppose the direction being tested.} \texttt{tab:rq1} reports one-sided $p$-values for the very hypothesis the paper sets out to establish.}{Use two-sided tests (the results will survive: the gaps are large), or explicitly justify the directional hypothesis from prior literature in Methods.}

\issue{B7}{\textbf{Fragmenter classification rests on 5 samples per question in small cells.} With $R=5$, ``no majority'' requires the mode $\le 2$ of 5; per-question consensus loss is a noisy Bernoulli measurement, and the HLE cells have $n=24$--$72$ questions. No uncertainty is reported on the consensus-loss percentages (e.g., ``Gemini 72.2\%'').}{Report Wilson CIs on the consensus-loss rates (question-level counts are available), and add a sentence on sensitivity of the majority threshold ($\ge 3/5$) --- e.g., how rates change under a $\ge 4/5$ definition.}

\issue{B8}{\textbf{GPQA question pools differ by model family, confounding cross-model comparison.} Claude models were run on a 79-question GPQA sample, GPT models on a 198-question sample. RQ1/RQ3 cross-model statements on GPQA (e.g., resister vs.\ capitulator groups) compare models on different items.}{Best: run the missing questions so all models share the 198-question pool (or restrict analysis to the intersection). If runs are too expensive, restrict all cross-model GPQA claims to within-family comparisons and flag the confound in both \S Datasets and Limitations, not just implicitly.}

\issue{B9}{\textbf{A slope of exactly $+0.00^\dagger$ is classified ``Capitulator.''} \texttt{tab:rq3}, GPT-5.4 Mini on GPQA: the group is assigned by the sign of a zero, non-significant slope in a cell that is at ceiling (98--100\% flip rates from $T_1$).}{Add a third label (``Ceiling'' or ``N.C.''---not classifiable) for cells with non-significant slope and near-ceiling baseline flip rate, and use it here.}

\issue{B10}{\textbf{The calibrator is never validated on Gemini traces.} Table \texttt{tab:cal\_overall} covers only the five Claude/GPT targets, yet RQ4/HLE conclusions (Table \texttt{tab:rq4\_hle}) include Gemini rows read by the same calibrator, and Gemini has the most extreme confidence values (0.978--0.990).}{Add a Gemini row to \texttt{tab:cal\_overall} and \texttt{tab:cal\_turn} (the predictions already exist), or scope the RQ4 claims to the validated five targets and mark the Gemini rows as unvalidated.}

\issue{B11}{\textbf{Experiment 5's ECE conflates target capitulation with calibrator error, but the main text presents it as model miscalibration.} The appendix (\S Turn-level calibration drift) correctly notes the two error sources cannot be separated, yet the Results paragraph ``Miscalibration under pressure'' attributes the ECE rise to the models (``Models do not become more confident under pressure; they become less accurate\ldots'').}{Carry the appendix caveat into the main text (one sentence), and where possible decompose: report the belief-accuracy drop (target capitulation) separately from the confidence trajectory, and attribute ECE growth to their divergence rather than to ``the model's expressed confidence'' alone. Also fix the notation: the ensemble formula reuses $K$ for what is actually the $R=5$ CoT runs (see D8).}

\end{issuelist}

\section{Claims vs.\ evidence and framing}
\label{sec:claims}
\begin{issuelist}

\issue{C1}{\textbf{Abstract overstates the three-group taxonomy as a model property.} ``Under escalating pressure models can be classified into three groups'' --- but the Discussion itself concludes the resister/capitulator label ``flips sign on HLE'' and is ``a joint property of (model, benchmark difficulty).''}{Rephrase the abstract: ``model--benchmark cells fall into three behavioral modes\ldots, and the mode is a joint property of model and benchmark difficulty rather than a fixed model trait.'' This is also the more interesting claim.}

\issue{C2}{\textbf{``Escalating pressure'' framing conflicts with the single-turn potency ordering.} The single-turn experiment shows category potency is non-monotonic in the sequence order (expert authority at $T_4$ strongest; crowd consensus at $T_6$ among the weakest), so the fixed sequence does not monotonically escalate persuasive force. The paper concedes this in a dense Methods aside, but the abstract, protocol name, and Table captions all say ``escalating.''}{Rename to ``fixed-order pressure sequence'' (or ``sustained multi-category pressure'') throughout, and promote the buried argument --- that monotone per-turn trajectories despite non-monotone category potency license the dose-accumulation reading of $\beta_P$ --- from Methods \S Statistical Models into the RQ3 results text, where readers need it.}

\issue{C3}{\textbf{The title promises prediction; the paper delivers association.} ``Uncertainty as a Predictor'' but no predictive analysis is reported (no AUC/ROC, no held-out prediction of flips from $H$ or $\hat h$).}{Cheapest strong fix: add a small predictive evaluation --- per-cell AUC of baseline $H$ (or $\hat h$) for predicting whether a dialogue ever flips. The data already exist and it would substantiate the title and the ``identify fragile questions before any adversarial interaction'' claim in the Conclusion. Alternatively, retitle (``\ldots Uncertainty Moderates Sycophantic Capitulation'').}

\issue{C4}{\textbf{Figure 1's caption overclaims uniformity.} It says uncertain questions flip ``at roughly twice the rate\ldots consistently across all five models''; for GPT-5.4 Nano the ratio is $76.1/65.4 \approx 1.16\times$ (though the difference is significant).}{Caption: ``flip at up to $7\times$ the rate (2$\times$ pooled), with a positive, significant gap for all five models.'' Keep ``all $p<0.001$'' pending the B1 re-analysis.}

\issue{C5}{\textbf{``Off-target drifts $<2\%$ of turns'' is stated globally but likely holds only on the easier benchmarks.} On HLE, five of six models lose consensus on most queries, implying substantial off-target mass.}{Scope the claim: ``fewer than 2\% of turns on MMLU-Pro and GPQA-Diamond; on HLE, off-target drift is the dominant mode and is analyzed separately as fragmentation (\S Results).'' Report the actual HLE off-target rate.}

\issue{C6}{\textbf{Deployment-safety claims outrun the evidence.} Discussion asserts ``Resisters are safest under extended adversarial dialogue,'' ``a confidence-threshold abstention rule\ldots would catch essentially none of the pressure-induced flips'' --- from six models, MCQ-only, one fixed prompt sequence, and (for the abstention claim) no abstention experiment.}{Hedge to what was measured: the abstention claim should either be tested (simulate a threshold rule on the existing $\hat\gamma$ trajectories --- the data exist) or phrased as an implication (``suggests\ldots would miss most flips, since confidence stays flat while belief collapses'').}

\issue{C7}{\textbf{Uncertainty vs.\ item-difficulty confound is handled quietly.} $H$ correlates with question hardness; a skeptic can read RQ1 as ``hard questions flip more.'' The $\hat h$ robustness check partially addresses this but is framed as a parameterization detail in the appendix.}{Add 2--3 sentences in Discussion distinguishing the behavioral-uncertainty reading from the item-difficulty reading, noting that (a) $\hat h$ is precisely calibrated hardness and gives the same story, and (b) the model-specificity of $H$ (same question, different cells) is what makes it a model-belief signal rather than an item property.}

\end{issuelist}

\section{Presentation and \LaTeX{} mechanics}
\label{sec:latex}
\begin{issuelist}

\issue{D1}{\textbf{Stray \texttt{\textbackslash centering} at top level.} \texttt{results.tex} lines 34--35: \texttt{\textbackslash centering} appears \emph{before} \texttt{\textbackslash begin\{figure\}}, so it leaks into the surrounding text block.}{Move \texttt{\textbackslash centering} inside the figure environment (and add the missing \texttt{[t]} placement to match sibling figures).}

\issue{D2}{\textbf{Figure 1 is never referenced.} \texttt{fig:overview} (the key-result teaser) has no \texttt{\textbackslash ref} anywhere; it also lives in a file misleadingly named \texttt{architecture\_diagram.tex}.}{Cite it from the Introduction (``\texttt{\textbackslash autoref\{fig:overview\}} previews the main result'') and rename the file \texttt{fig1\_overview.tex}.}

\issue{D3}{\textbf{Dead files and dead tables in the source tree.} Not \texttt{\textbackslash input} anywhere: \texttt{tables/rq2.tex}, \texttt{tables/rq2\_hardness.tex} (superseded by \texttt{rq2\_combined.tex}), \texttt{figures/arch\_static.tex}, \texttt{figures/arch\_with\_adversarial.tex}, \texttt{figures/worked\_example.tex}, \texttt{sections/fig/rq2\_coefs\_v1.png}, \texttt{sections/fig/rq2\_interaction\_combined.png}. Duplicate labels lurk if they are ever re-included.}{Delete them from the submission (keep in version control if needed).}

\issue{D4}{\textbf{Commented-out drafts and internal review notes remain in the source.} Old abstract (\texttt{abstract.tex}), old conclusion (mentions ``corrigibility gap'' and ``Constitutional AI,'' concepts absent from the paper), old limitations, a $\sim$300-line commented table in \texttt{tables/rq4.tex}, and --- most important --- a commented Discussion paragraph containing the note ``\emph{the evidence does not support this claim. There is no significant positive interaction},'' with author names in the todo macros. arXiv publishes source; this is a real leak risk and de-anonymizes.}{Strip all commented-out blocks and the \texttt{todonotes} author macros before submission (\texttt{latexpand --empty-comments} or arXiv's own comment stripping; verify manually).}

\issue{D5}{\textbf{Nonstandard negated-entropy convention imposes a tax on every section.} Defining $H=\sum \hat p \log \hat p \le 0$ forces sign-convention reminders in Methods (twice), Results, and three table captions, and produces double-negative prose (``more negative = more spread'').}{Use standard Shannon entropy $H \ge 0$ and flip the signs of $\beta_H$, $B_t$, and the axis conventions. Every ``sign convention'' sentence can then be deleted, and $\beta_H > 0$ reads naturally as ``more uncertainty, more flips.''}

\issue{D6}{\textbf{Missing cross-reference glue.} \texttt{methods.tex}: ``\ldots on their own \texttt{\textbackslash autoref\{tab:single\_turn\}}'' and ``\ldots sampling parameters \texttt{\textbackslash autoref\{app:hyperparams\}}'' render as bare refs jammed into sentences.}{Wrap as ``(\texttt{\textbackslash autoref\{...\}})'' or ``see \texttt{\textbackslash autoref\{...\}}''; do a global pass for consistent \texttt{\textbackslash autoref} vs.\ \texttt{Table\textasciitilde\textbackslash ref} usage.}

\issue{D7}{\textbf{Duplicate \texttt{\textbackslash appendix}.} Issued in \texttt{main.tex} (line 97) and again in \texttt{sections/appendix.tex} (line 2).}{Keep exactly one (in \texttt{main.tex}); remove the other plus the redundant \texttt{\textbackslash clearpage}.}

\issue{D8}{\textbf{Notation collisions.} $K$ means 25 baseline samples in \S3.2 but is reused for the $R=5$ CoT runs in Experiment 5's ensemble-confidence formula; the appendix mixed-model equation uses $H_{m,q}$ where Methods uses $H_{c,m,q}$; $\hat p_{q,t}$ in the Flip Definition lacks the dialogue index $r$ used two paragraphs earlier.}{Rename Experiment 5's index to $r=1..R$; make the appendix GLMM formula use $H_c$; add the $r$ subscript in the flip definition.}

\issue{D9}{\textbf{Small mechanical nits.} \texttt{tab:rq3} caption: ``ClaudeHaiku'' (missing space); ``Mann-Whitney'' should be ``Mann--Whitney'' (en-dash) in \texttt{tab:rq1}; ``GPQA-Diam.''\ vs.\ ``GPQA-Diamond'' used inconsistently; \texttt{todonotes} margin-width warning; two \texttt{[h]} floats coerced to \texttt{[ht]}; missing emails for two co-authors in \texttt{main.tex}.}{Fix strings; load \texttt{todonotes} with \texttt{[disable]} (or drop it per D4); use \texttt{[t]} float specifiers; complete author metadata before the camera-ready.}

\issue{D10}{\textbf{RQ1 caption will be invalidated by the A3 fix.} ``N/A indicates a cell with no certain group (every question is uncertain at baseline)'' is tied to the contested claim.}{Update the caption together with A3 so caption, text, and appendix agree.}

\end{issuelist}

\section{Reproducibility and TMLR compliance}
\label{sec:tmlr}
\begin{issuelist}

\issue{E1}{\textbf{All 35 bibliography entries are arXiv \texttt{@misc} entries.} Several are published: Sharma et al.\ (ICLR 2024), Turpin et al.\ (NeurIPS 2023), Kuhn et al.\ (ICLR 2023), Tian et al.\ (EMNLP 2023), Perez et al.\ (Findings of ACL 2023), MMLU-Pro (NeurIPS 2024 D\&B), GPQA (COLM 2024), Wei et al., etc.}{Replace with the published venue entries (DBLP export), keeping arXiv IDs as \texttt{eprint} fields.}

\issue{E2}{\textbf{No code/data availability statement, but the appendix references ``the released code.''} Table \texttt{tab:hyperparams} says ``Defaults are from the released code'' with no pointer; TMLR reviewers weigh reproducibility heavily.}{Add an anonymized repository (transcripts, per-question $H$, regression code) and a Reproducibility Statement section; TMLR explicitly invites one.}

\issue{E3}{\textbf{Sampling procedures for benchmark subsets are unstated.} How were the 420 MMLU-Pro questions, the 79- vs.\ 198-question GPQA samples, and the 240 HLE questions selected (random seed, category balance, MC filtering rule)?}{Add a short ``Subset construction'' paragraph in \S Datasets with the sampling rule and seed for each; release the question ID lists with E2.}

\issue{E4}{\textbf{No Broader Impact Statement.} The paper documents which rhetorical strategies most effectively flip models (expert authority) --- mildly dual-use knowledge --- and has direct deployment implications.}{Add TMLR's optional Broader Impact Statement: one paragraph on the defensive intent, the responsible framing of the potency ordering, and the deployment guidance.}

\issue{E5}{\textbf{API model snapshots are moving targets.} Snapshot names and the May--June 2026 window are given (good), but results cannot be regenerated once providers rotate defaults.}{Alongside E2, release the raw transcripts so all analyses are re-runnable without API access; state this explicitly.}

\end{issuelist}

\section{Writing and organization}
\label{sec:writing}
\begin{issuelist}

\issue{F1}{\textbf{Introduction's final sentence is a run-on with three unrelated payloads.} ``We evaluate five OpenAI and Anthropic models\ldots and add Gemini 3.5 Flash on Humanity's Last Exam\ldots to understand the effect of uncertainty and hardness on the eventual flip\ldots and if the eventual flip could be identified through the reasoning trace\ldots''}{Split into (1) evaluation scope, (2) what RQ1--RQ3 establish, (3) what RQ4 adds. Mirror the RQ list.}

\issue{F2}{\textbf{Abstract's coverage parenthetical is convoluted.} ``\ldots and extend to HLE for all six targets (adding Gemini 3.5 Flash on the HLE stress test and single-turn category comparison, where its eligibility criterion is met).''}{Simplify: ``We evaluate six GPT, Claude, and Gemini models on MMLU-Pro, GPQA-Diamond, and HLE; cells without baseline uncertainty variance are excluded from uncertainty-conditioned analyses (\S X).'' Details belong in Methods.}

\issue{F3}{\textbf{Grammar in Background.} ``Prior work \texttt{\textbackslash citet\{xu2023earthisflat\}} show that\ldots'' (subject mismatch), and ``\ldots study that models often fail\ldots'' should be ``find that'' or ``show that.''}{Rewrite as ``\texttt{\textbackslash citet} show that\ldots'' / ``\texttt{\textbackslash citet} find that models often fail\ldots''; proofread the section for verb agreement.}

\issue{F4}{\textbf{Experiment numbering drifts from the RQ structure.} Experiments 1--4 map to RQ1--RQ4; then an unnumbered ``Single-turn robustness check'' paragraph and an ``Experiment 5: Ensemble calibration'' appear, the latter answering no stated RQ.}{Fold Experiment 5 into Experiment 4 (it quantifies the same dissociation) or introduce it as a named analysis under RQ4; give the single-turn check a consistent label used in Results (``Single-turn baseline'') in both places.}

\issue{F5}{\textbf{Methods \S3.1 (signal overview) substantially duplicates \S3.2--\S3.5.} The four-signal reference list re-explains definitions that are then formally re-introduced, adding $\sim$1 column of length.}{Compress \S3.1 to a table (signal, granularity, defined-in, used-for) plus the sign-convention note (or delete the latter after D5).}

\issue{F6}{\textbf{Key methodological argument is buried in a caption-length Methods aside.} The mean-centering rationale, the separation caveat, and the accumulation-vs.-potency argument all live inside one dense paragraph in \S Statistical Models.}{Split into three short labeled paragraphs (``Centering,'' ``Separation,'' ``Order confound''); move the order-confound argument's conclusion to Results (see C2).}

\end{issuelist}

\vspace{1em}
\progress